\BeleskeID{09_2-s021}
% element: 09_2-s021:e05
Polazna ravnoteža struja u oblasti niskog izlaznog napona je
\[k_p(-V_{DD}-V_{Tp})^2=k_nV_O\left(2(V_I-V_{Tn})-V_O\right)\]
Za ulaznu logičku jedinicu dobijamo jednačinu prikazanu na slajdu. Ako je $V_{OL}$ mali, zanemaruje se njegov kvadrat:
\[k_p(-V_{DD}-V_{Tp})^2\approx k_nV_{OL}\left(2(V_{OH}-V_{Tn})\right)\]
Dobijeni približan rezultat sadrži odnos parametara $k_p/k_n$. Uz jednake apsolutne pragove, $V_{Tp}=-V_{Tn}$, jedan faktor $V_{DD}-V_{Tn}$ se skrati i izraz se svodi na poslednji oblik. Pretpostavka malog izlaznog napona mora se proveriti sa izabranim odnosom parametara.
